% Fresh source partition and proof-use audit, 2026-10-22.
\section{Every canonical C2 source: an exhaustive mathematical join}
\label{sec:c2-complete-source-routing}

\paragraph{Scope and notation.}
Let \(K_c,K_t\) be arbitrarily chosen, separately labelled real closed
extensions of \(\mathbb R\); no algebraic-disjointness assumption is made.
All source inequalities below are in \(K_c\).
Use the entire original physical C2 parent
\[
\begin{split}
 F={}&H+4nD_uD_vX^2+4mM^2+4CD_v(uW-Y)^2
       +4AD_uD_vY^2+4DJ^2+4\alpha D_u(vY-Z)^2,\\
 M={}&vW-X+uvY-uZ,\qquad J=uW+uvX-Y-vZ,\\
 H/2={}&J^2+D_u(vY-Z)^2+D_v(X+uZ)^2-D_uD_vZ^2 .
\end{split}\tag{CT1}
\]
Here \(D_u=1+u^2,D_v=1+v^2\), \(d=1\), and the omitted weights are
\(\lambda_{\rm event}=\beta=0\).  Assume an actual finite noncommuting
stationary zero.  The native-entry theorem EN gives exactly
\[
 m,C,A,D,\alpha>0,\quad n\geq0,\quad
 c,a,b,t>0,\quad 0<r<1,\quad
 h=b-a^2t>0,\quad g>0,\quad j>0,
\]
\[
 g=(1+c^2)b-(1-r+a^2c^2t),\qquad j=1+c^2r-t,\qquad
 1-c^2h>0,\quad b-a^2c^2r>0.                    \tag{CT2}
\]
The actual angles are \(u_c=c,v_c=ac\), and the six coefficients
are the actual rational stationary coefficients, with the nonzero
stationary determinant proved in EN.  In particular no limiting or
formally perturbed source is substituted for CT1.

\paragraph{Theorem.}
Under CT1--CT2, the full parent is a finite sum of squares in
\(K_c\otimes_{\mathbb R}K_t\) at every projective testing pair and
every real or complex state over \(K_t\).
All coefficient scales, literal \(n=0\), and every positive source
opening allowed by CT2 are included.  This theorem has no
earlier-admission-complement hypothesis.

\subsection{Two finite partitions}

First put
\[
 \delta_{\rm tot}=(1+n+m+C+A+D+\alpha)^{-1}.
\tag{CT3}
\]
If \(\delta_{\rm tot}\) is a unit, all six coefficients are bounded,
and the complete positive-scale theorem AU applies at its now-proved
native-entry interface.  Otherwise \(\delta_{\rm tot}\) is positive
infinitesimal.  No normalization of an individual large event is
used to replace this dichotomy.

Call a positive event weak when it is infinitesimal after this
normalization \(d=1\), and strong otherwise.  A strong event may be
unbounded.  If exactly one of \(C,\alpha\) is weak, orient it to be
\(C\), using the exact full-parent reflection
\[
\begin{split}
 &(u,v;W,X,Y,Z;n,m,C,A,D,\alpha)\\
 &\quad\longmapsto
 (1/v,1/u;Z,X,Y,W;n,D,\alpha,A,m,C).
\end{split}\tag{CT4}
\]
In homogeneous angle coordinates the identity is
\(F=u^2v^2F'\); its polynomial homogeneous interpretation includes
literal testing zero and infinity.  CT4 preserves all source gaps:
writing \(z=c^2\), its parameters are
\[
 c'=(ac)^{-1},\quad a'=a,\quad b'=a^2/b,\quad
 r'=a^2zr/b,\quad t'=(1-zh)/b,
\]
\[
 h'=a^2zh/b,\quad g'=a^2j/b,\quad j'=g/b,\qquad
 \Delta'=\Delta/(b^2z^2),\quad N'=a^2N/b^2 .
\tag{CT5}
\]
These are pure-source identities, and give positive values for
\(1-r',1-z'h',b'-a^2z'r'\).  They preserve
\(\delta_{\rm tot}\) exactly.  EN and HJ verify the full reflection;
the retained matching-source entry prints the displayed rational
source identities.  No assertion \(t\mapsto1/t\) is used.

Let \({\rm I},{\rm U},{\rm O}\) denote, respectively, a positive
infinitesimal, a positive bounded unit, and a positive unbounded
source number.  Every positive actual angle has exactly one of
these types.  After the two-weak case is treated below, the
remaining dispatch is
\[
\begin{array}{c|c|c}
(u_c,v_c)&C,\alpha\text{ both strong}
                 &C\text{ weak},\ \alpha\text{ strong}\\ \hline
{\rm II}&\text{strong II local and remote}&\text{impossible}\\
{\rm IU}&\text{single-angle join}&\text{impossible}\\
{\rm IO}&\text{strong IH local and remote}&\text{weak IH full join}\\
{\rm UI}&\text{strong UI full join}&\text{impossible}\\
{\rm UU}&\text{five retained events}&\text{impossible}\\
{\rm UO}&\text{single-angle join by CT4}&\text{single-angle join by CT4}\\
{\rm OI}&\text{strong HI full join}&\text{impossible}\\
{\rm OU}&\text{strong UI full join by CT4}&\text{impossible}\\
{\rm OO}&\text{strong II full join by CT4}&\text{HJ by CT4}
\end{array}\tag{CT6}
\]
The letters H in the retained file titles mean the type O in this
table.  They denote an unbounded \emph{finite} actual half-angle,
not a literal projective endpoint.  Testing endpoints are included
by each full-parent theorem.

\subsection{The two-weak case and the empty cells}

Suppose first that \(C,\alpha\) are weak.  The actual proof
\emph{chart\_free\_c2.tex} supplies the following source implications,
without a tensor-SOS premise.  The coefficient-field experiments
which kill three matching packets show that
\[
                    n,m,A,D\text{ are unbounded}.             \tag{CT7}
\]
For example the remaining \(D\) experiment has full value
\(4D\varepsilon^2-2\varepsilon+4(2+\varepsilon)C
+8(1+\varepsilon)\alpha\).
If \(D\) were bounded, a sufficiently small positive real
\(\varepsilon\) would contradict physicality.  The other three
experiments are its exact event-cycle copies.
The exact signed-core estimate
\[
 H\geq-2u_cv_c\|\psi_c\|^2
\]
follows from the printed matrix-square identity
\((M_H-LI)^2=L^2I+4u_cv_cLJ_0\), \(J_0^2=I\).
Each matching amplitude divided by \(\sqrt{u_cv_c}\) is therefore
infinitesimal at a norm-one contact.  Compact angle coordinates
and, if needed, CT4 imply
\[
 c\text{ bounded},\qquad a,b,t\sim1,\qquad
 r,h,g,j\text{ positive infinitesimal}.                    \tag{CT8}
\]
In the small-angle case the two matching equations divided by
\(\sqrt{u_cv_c}\) first force \(v_c/u_c\sim1\) and \(W_c/Z_c\sim1\);
this is a relative bound, not just a compact limiting statement.
The same body proves \(\alpha/u_c^2\) infinitesimal before using
scaled stationarity.  Its invertible four-amplitude Jacobian has
residue determinant \(2\), and the weighted stationary gradients
give \(n\ell_n,A\ell_A,m\ell_m,D\ell_D\) residues
\((-1,1,1,-1)/4\).  Thus all strict signs in CT8 are justified.

For \(c\) infinitesimal use the complete six-root radial construction
in \emph{c2\_unified.tex}, \emph{c2\_y\_current.tex}, and
\emph{c2\_radial\_corner.tex}.  Only its radial domain CT8 is imported.
Its pure-source weight \(N\) satisfies
\[
 N\asymp r+h+g+j,\qquad
 r/N,h/N,g/N,j/N=O(1),\qquad
 \operatorname{res}\bigl(4\alpha/(c^4N)\bigr)=1.          \tag{CT9}
\]
The corrected radial root is \(\sqrt N(u^2-c^2)\).
The actual proof retains all four normalized event roots and every
state normal; the final exact current-zero correction makes every
radial cross coefficient belong to \((r,h,g,j)\) before division
by \(\sqrt N\).  Its Gram has strict residue on the full six-root
list, including the complementary \(W\)-state chart.
The actual outer proof \emph{c2\_radial\_all\_tests.tex}
uses
\[
             1/(Cn),\ 1/(Cm),\ 1/(CA),\ 1/(CD)=O(1),
\]
which follow from
CT9 and the unchanged source formulas.  Its reciprocal-corner
conductor and its two complete state charts retain the weak \(C\)
packet exactly.
For \(c\) a unit, \emph{c2\_matching\_strength.tex} gives
\(\alpha/C\asymp1\), both sets of four product bounds, and the
same outer proof after CT4 covers the remaining corner.
Both cases use the detailed arbitrary-scale physical-source
noncommuting-residue theorem A6 for the noncommuting testing
complement.  A6's full geometry, curvature, weighted matching,
promotion, and induction bodies are explicit dependencies.
This completes every two-weak source, rather than assuming a
preselected radial chart.

For the right column of CT6, the exact proof
\emph{angle\_reductions.tex} establishes the II, IU, UI and UU
contradictions from CT2 and \(n\geq0\).
Some useful details fix its scope.  On II, putting
\[
 H=c^2h,\quad L=c^2(b-h),\quad B=H+L
\]
in its own actual-angle notation makes the positive numerator
bound force \(\alpha\) weak whenever \(C\) is weak.  On UI, the
same source numerator and the \(n\geq0\) inequality force
\(r\) infinitesimal and \(h\) a unit, contradicting weakness
of \(C\).  On IU, bounded \(b\) is a unit by the inequality in
CT10 below and gives \(C\) strong; unbounded \(b\) gives
\(C\asymp L^{-1}\) with \(L\) bounded.  On UU, \(a,c\) are units,
\(\Delta\) and \(b+r\) are units, and the positive numerator
identity makes \(C,\alpha\) weak together or strong together.
These are necessary-source arguments, not all-tests conclusions.

The OI and OU contradictions use
\emph{bounded\_second\_angle\_exclusion.tex}.
With \(q=u_c^{-2}\) infinitesimal and \(y=v_c^2\) bounded, its
exact numerator for \(n\) is \(-qP/(2rQ)\), \(Q>0\).
The printed positive remainder gives
\[
 P\geq r^2(1+y)(H+yr)-Tqr(y+2)-q^2(T^2+T),\qquad
 0<T<H+y .
\]
If \(q/r\) were infinitesimal, division by \(r^2(H+yr)\)
would force \(P>0\), contradicting \(n\geq0\).
Thus \(r/q\) is bounded, \(H\sim1\), and the exact cancelled
\(C\) formula has a positive-unit numerator and bounded
positive denominator.  Consequently \(C\) is strong.  Literal
\(n=0\) obeys the same contradiction.

\subsection{The single-angle join has no omitted opening}

For an IU source put \(x=u_c^2\) infinitesimal and
\(k=v_c\) a unit.  If \(b\) is bounded, the exact source
inequality \(n\geq0\) gives
\[
 r_0\{b_0+k_0^2(b_0+1)\}\leq b_0,\qquad
 r_0\geq1-b_0,\qquad
 b_0\geq k_0/\sqrt{1+k_0^2}>0 .                       \tag{CT10}
\]
The actual body \emph{c2\_bounded\_state\_entry.tex} proves
the last bound also when \(b_0\geq1\).
There are then exactly three source alternatives:
\[
\begin{array}{c|l}
g\text{ unit}&\text{three-opening proof, arbitrary }r,h,b-h>0,\\
g\text{ infinitesimal},\ r\text{ unit}
             &\text{last-opening proof},\\
g,r\text{ infinitesimal}
             &\text{joint centered-packet proof}.
\end{array}\tag{CT11}
\]
All bounded positive \(g,r\) have been included.
The first proof supplies both bounded and unbounded \(n\)
normalizations; the second uses the triangular hard normal of
the entire \(m\) amplitude; the third retains both contact-zero
centered \(n,m\) packets.  Their common full scalar proof is the
six-term balanced Laurent decomposition in
\emph{general\_single\_angle.tex}, with its five-term limiting
version and joint version proved in the respective bodies.
This supplies global equality uniqueness and a positive contact
Hessian.  The same bodies print complete first-normal divisions,
hard blocks, the literal \(n=0\) double-zero corner and all four
outer-core identities.  Their remaining testing complement uses
the actual detailed A6 theorem.  No normal is omitted and no
bound on \(h\) or \(b-h\) is restored implicitly.

If \(b\) is unbounded, put \(B=xb,H=xh,L=x(b-h)\).
The original gaps give
\[
 B=H+L,\quad 0<H<1,\quad 0<L<k^2(1+xr).
\]
Thus \(B\) is bounded.  The exact proof
\emph{unbounded\_b\_reduction.tex} cancels the original
denominators and gives
\[
 1-r\asymp1,\quad m\asymp b^{-2},\quad
 C\asymp L^{-1},\quad A\asymp H^{-1},\quad
 D\asymp B/j,\quad \alpha\asymp B/(1-H).               \tag{CT12}
\]
If \(B\) is a unit, \(C,\alpha\) are strong and the full strong
pitchfork below applies.  If \(B\) is infinitesimal, the actual
\emph{unbounded\_b\_small\_B.tex} theorem applies with
\(\sigma=1/b\), with no restriction on \(B/\sigma\).
Its separate homogeneous scalar lemma proves the entire
contact-ideal lift on
\[
 (uW-(B/k)Z)/\sqrt B,\quad uZ-\sigma kW,\quad
 \sqrt B(v-k)Z,\quad u(v-k)Z .
\tag{CT13}
\]
The whole-parent proof retains the possibly nonzero first-normal
angular coupling and buys its strict pivot by the exact zero
\[
 T^2-\sigma U^2-ST(v-k)-\sigma kRU(v-k)=0 .
\]
It proves both full normals, arbitrary \(W/Z\), every \(n\) scale,
all openings, and the literal \(n=0\) corner.  Its outer charts
use the displayed general-single-angle identities with scale
\(\sqrt B\), and the detailed A6 theorem for the final complement.
This exhausts IU.  CT4 exhausts UO with no new source assumption.

The strong pitchfork invoked in CT12 is itself complete at all
upper scales.  Its three-large-event body covers bounded-unit
\(C,\alpha\), with \(n,A,D\) arbitrary; its two-large-event
predecessor supplies the printed stationary equations, full
normal Taylor division and real boundary Schur calculation.
The body \emph{strong\_complement\_pitchfork.tex} treats
unbounded \(C\) or \(\alpha\).  Its coefficient-field state
equations force \(m/k^4\asymp1\) and bounded weighted graph
coordinates before any small gap is inverted.  It uses the
actual large \(n,C,\alpha\) or \(n,D,\alpha\) packets as normals.
The five-root full-parent Gram has a strict angular Schur pivot
because the normalized real boundary numerator is
\(q_{12}(v^2-a_0^2)^2\) with \(q_{12}>0\).
Only this real residue Schur computation uses a determinant
inverse; no mixed tensor denominator is inverted.
For its outer charts one may use the complete cap and corner
proof printed in that same strong-complement body, replacing
the older predecessor's reference to a preceding outer corollary.

\subsection{The remaining occupied cells}

\paragraph{II and OO with both complementary events strong.}
The exact common-II source reduction proves that
\(\alpha\) is strong precisely when \(B=H+L\) is a unit.
Then \(L\) is infinitesimal, \(H\) and \(1-r\) are units, and
\[
 m/z^2,\ A/y,\ CL,\ Dyj,\ \alpha(1-H)\asymp1,
 \qquad nr=O(1).                                      \tag{CT14}
\]
The exact numerator at \(r=1\) is negative, giving the asserted
unit gap \(1-r\); none of \(r,j,1-H\) is assumed a unit.
The proof \emph{local\_strong\_ii.tex} retains all four centered
large-packet excesses and gives the full three-normal remainder.
Its angular identity, harmonic \(C,D\) square and free radial
pivot are exact; the normalized hard residue and the remaining
two-by-two Schur margin are strictly positive.
The separate \emph{strong\_ii\_remote.tex} uses the original
homogeneous parent and harmonic current, not a higher-degree
homogenization of the local graph.  It proves all remaining
angles and state charts, including the endpoint forms.
Thus there is no local-to-remote scope gap.  CT4 gives OO.

\paragraph{IO, weak or strong \(C\).}
The weak IH source entry proves \(B\) unbounded and
\(\rho=B/y\) bounded with \(0\leq\bar\rho\leq1\).
The full interior-corner proof, both \(\rho=0,1\) source-endpoint
extensions, remote-chart proof, and separate open-axis proof
cover all independent angles and states.  The latter open-axis
proof is indispensable: it supplies the product-pivot split
when \(u\) is a unit and \(1/v\) is infinitesimal.
All these bodies retain \(nr\) and the full \(D\) packet at the
closing endpoint; literal \(n=0\) is included.
The actual weak-IH all-tests join merely assembles these
already specified bodies.

If \(C,\alpha\) are both strong, the separate strong-angle
source entry gives \(B\) a bounded unit, \(r=O(1/y)\),
\(Dy^2\asymp1\) and the remaining opening-safe products.
The strong-IH local proof reuses only the explicit polynomial
current identities of the weak-IH proof, with its own source
estimates and finite \(C\) cap; it does not import the weak-IH
source conclusion.  Its remote proof includes the two open
axes and reciprocal corner and uses detailed A6 only when
both testing pairs are noncommuting in residue.

\paragraph{UI and OU with both complementary events strong.}
The actual \emph{full\_parent\_reduction.tex} in the C2-UI
directory cancels all closing source openings before dividing
the normals.  In particular \(nr,m\mu,\alpha\pi,C\rho k^2,
D\sigma k^2\), and the difference \(Ct-Dj\), are bounded.
The exact state inverse is polynomial and never divides by
a state coordinate.  The reciprocal-contact lemma and local
certificate prove the full contact Gram, including all three
normals.  The complete \emph{ui\_all\_tests\_recovered.tex}
proof supplies real interior-source approximation for its
finite-cap step and strict Grams on every complementary
axis and state chart.  Its Laurent equality and positive
contact Hessian follow from the actual source equations.
No new global source classification is imported from this
local theorem.  CT4 gives OU.

\paragraph{OI.}
The exact HI source entry proves directly from \(n\geq0\)
that \(R=r/q\) is bounded, \(H\sim1\), and
\[
 mqK,\ Cyt,\ Dyj,\ \alpha qE,\ A/(qy)\asymp1,
 \quad q=u_c^{-2},\ y=v_c^2\text{ infinitesimal}.       \tag{CT15}
\]
Thus \(m,C,D,\alpha\) are all unbounded, without a prior
strong-event assumption.  Its remote full-parent theorem
caps these four coefficients at a sufficiently large
\emph{real} constant after fixing the testing residue.
It prints strict full state forms on every remote chart
and literal critical axis; only \(1/u,v\) both nonzero
infinitesimal remain.
The retained-contact module gives the complete local
four-monomial circuit and radical conductor there.
The full HI module includes the actual \(n\) packet, its
two complete normals, and the endpoint \(R_0=0\) where
\(n\) can be unbounded.  Its constant circuit multiplier
lies in \([2/3,2]\), and its normalized hard residue
is strictly positive.  Its proof uses the predecessor
only for the explicit half-angle conductor, never for
an unsupported \(n=0\)-to-\(n>0\) inference.
Together these three bodies close OI.

\paragraph{UU and the five-event terminal.}
The exact source-angle equations give \(m,A,D\) strong on
the both-strong UU cell.  The actual five-retained-events
proof uses \(\mu=\min(1,m,C,A,D,\alpha)\), a positive unit,
and retains all excess original packets.
Its commuting-axis full state forms and three nontrivial
double-corner coordinate Grams are strict.  For
noncommuting testing pairs it uses the detailed arbitrary-scale
physical-source theorem A6 for the original source.
The selected capped parent is never asserted globally
physical.  This distinction prevents a contact-zero
packet from being improperly subtracted at the source
contact.

\paragraph{OO in the directed weak case.}
CT4 gives exactly the native II domain HJ1:
\(C\) strong, \(\alpha\) weak, both actual angles
infinitesimal, and \(\delta_{\rm tot}\) infinitesimal.
HJ1--HJ16 prove the unchanged source equations, the
equivalence of the weak face with \(B\) infinitesimal,
the complete ten-row partition, all its literal \(n=0\)
maps, and every source/testing/state join.
The bare-source audit BSR1--BSR15 proves the precise
retained R5--R10 hypotheses and restores each actual
earlier exit.  Neither HJ nor BSR assumes this C2 theorem,
the whole C1 theorem, or AU.  Its only C1 uses are the
two explicit bare CS and small-e all-tests bodies under
the exact \(n=0\) source map.
Consequently this last cell introduces no dependency cycle.

\subsection{Proof-use scope and the conclusion}

The mathematical input list accompanying this theorem pins
every actual proof body used above, including the subordinate
single-angle and homogeneous scalar proofs, the complete normal
and source-endpoint constructions, and the noncommuting residue proofs.
An inherited final paragraph which describes a then-open
complement is not a premise.  Only the displayed source
hypotheses and the exact full-parent constructions just
specified are imported.
In particular:
\[
\begin{gathered}
 \text{native entry is EN, not an unavailable sign label;}\\
 \text{positive scale is AU, not an undefined terminal call;}\\
 \text{weak II is HJ, not its former frontier assertion.}
\end{gathered}\tag{CT16}
\]
The bodies in CT10--CT15 use the exact source equations,
positivity, the typed source scales and the displayed
opening bounds.  Their scalar divisions, normal bounds,
caps and outer joins do not use exclusion from any
blanket list of earlier admissions.
Those exclusions are therefore absent from the theorem:
this is a proof-use verification on the actual bodies,
not deletion of a logical hypothesis by relabelling.

Every positive \(\delta_{\rm tot}\) has one of the two
types in CT3.  In the infinitesimal case \(C,\alpha\)
have exactly one of the three weak/strong patterns
used above; CT4 orients the asymmetric pattern.
The nine ordered angle types in CT6 are exhaustive.
Every occupied cell has a full-parent theorem at its
exact source hypotheses, and every empty cell has the
necessary-source contradiction stated above.
Thus CT1 is tensor SOS at every independent test.
Complex states split into real and imaginary parts
because CT1 is a real quadratic form in the state.
All source and testing normalizations are pure-field
square factors.  Restoring the original positive \(d\)
completes the theorem.
